% Options for packages loaded elsewhere
\PassOptionsToPackage{unicode}{hyperref}
\PassOptionsToPackage{hyphens}{url}
%
\documentclass[
  10pt,
  a4paper,
]{article}
\usepackage{amsmath,amssymb}
\usepackage{iftex}
\ifPDFTeX
  \usepackage[T1]{fontenc}
  \usepackage[utf8]{inputenc}
  \usepackage{textcomp} % provide euro and other symbols
\else % if luatex or xetex
  \usepackage{unicode-math} % this also loads fontspec
  \defaultfontfeatures{Scale=MatchLowercase}
  \defaultfontfeatures[\rmfamily]{Ligatures=TeX,Scale=1}
\fi
\usepackage{lmodern}
\ifPDFTeX\else
  % xetex/luatex font selection
  \setmainfont[]{DejaVu Serif}
  \setmonofont[]{DejaVu Sans Mono}
  \setmathfont[]{Latin Modern Math}
\fi
% Use upquote if available, for straight quotes in verbatim environments
\IfFileExists{upquote.sty}{\usepackage{upquote}}{}
\IfFileExists{microtype.sty}{% use microtype if available
  \usepackage[]{microtype}
  \UseMicrotypeSet[protrusion]{basicmath} % disable protrusion for tt fonts
}{}
\makeatletter
\@ifundefined{KOMAClassName}{% if non-KOMA class
  \IfFileExists{parskip.sty}{%
    \usepackage{parskip}
  }{% else
    \setlength{\parindent}{0pt}
    \setlength{\parskip}{6pt plus 2pt minus 1pt}}
}{% if KOMA class
  \KOMAoptions{parskip=half}}
\makeatother
\usepackage{xcolor}
\usepackage[margin=20mm]{geometry}
\setlength{\emergencystretch}{3em} % prevent overfull lines
\providecommand{\tightlist}{%
  \setlength{\itemsep}{0pt}\setlength{\parskip}{0pt}}
\setcounter{secnumdepth}{-\maxdimen} % remove section numbering
\ifLuaTeX
\usepackage[bidi=basic]{babel}
\else
\usepackage[bidi=default]{babel}
\fi
\babelprovide[main,import]{english}
\ifPDFTeX
\else
\babelfont{rm}[]{DejaVu Serif}
\fi
% get rid of language-specific shorthands (see #6817):
\let\LanguageShortHands\languageshorthands
\def\languageshorthands#1{}
\usepackage{amsmath,amssymb,amsthm,mathtools,graphicx,pdfpages,booktabs,longtable,needspace}
\usepackage{fvextra}
\DefineVerbatimEnvironment{Highlighting}{Verbatim}{commandchars=\\\{\},breaklines=true,breakanywhere=true,fontsize=\small}
\setlength{\emergencystretch}{3em}
\allowdisplaybreaks
\raggedbottom\sloppy
\widowpenalty=10000\clubpenalty=10000
\providecommand{\Tr}{\operatorname{Tr}}
\providecommand{\tr}{\operatorname{Tr}}
\ifLuaTeX
  \usepackage{selnolig}  % disable illegal ligatures
\fi
\usepackage{bookmark}
\IfFileExists{xurl.sty}{\usepackage{xurl}}{} % add URL line breaks if available
\urlstyle{same}
\hypersetup{
  pdftitle={Fourth-Moment Control of the Wilson Nonlinear Correction},
  pdfauthor={Kirtus Justus (Theoretical Physicist)  31 October 2026Port Coquitlam, British Columbia, Canada},
  pdflang={en},
  hidelinks,
  pdfcreator={LaTeX via pandoc}}

\title{Fourth-Moment Control of the Wilson Nonlinear Correction}
\author{Kirtus Justus (Theoretical Physicist) \textbullet{} 31 October 2026\tabularnewline Port Coquitlam, British Columbia, Canada}
\date{}

\begin{document}
\maketitle

\subsection{Abstract}\label{abstract}

For the mass-scaled Wilson--Dirac interval covariance, we derive a vanishing bound for the transverse nonlinear correction, sharpening a square-root-logarithmic upper bound. The argument combines the proven positive spectral-majorant tail estimate with the sine-kernel/Jacobi second moment. Specifically, the transverse contribution is \(O_\mu(\sqrt{\log N/N})\), while the longitudinal contribution remains \(O_\mu(1)\). Consequently the complete nonlinear correction is uniformly bounded at fixed scaled mass. Together with a separately established cross-boundary entropy bound, this reduces the leading relative-entropy problem to its exact reference-modular residual, without proving the latter vanishes or has a prescribed asymptotic order.

\subsection{1. Model and inputs}\label{model-and-inputs}

Let \(N=2L\) and \(m=\mu/L\) with fixed \(\mu>0\). Consider the interval compression of the occupied-band projection of

\[
h_m(k)=\sin k\,\sigma_x+(1+m-\cos k)\sigma_z.
\]

Under the established massless scalarization, the reference covariance is \(B_N\), with commuting Jacobi matrix \(J_N\). Write \(W_N=B_N(I-B_N)\), and decompose the perturbation as \(\Delta_N=S_N+D_N\). Let \(v_j\) be common eigenvectors of \(B_N\) and \(J_N\), with eigenvalues \(b_j\in(0,1)\) and \(\lambda_j\); put \(q_j=\langle v_j,G_Nv_j\rangle\ge0\).

The already proved Wilson spectral-majorant theorem supplies

\[
\lvert\langle v_j,S_Nv_j\rangle\rvert\le q_j,\qquad
\lvert\langle v_j,D_Nv_j\rangle\rvert\le2\sqrt{q_jb_j(1-b_j)},
\] \[
\sum_{|\lambda_j|>t}q_j\le\frac{a_\mu}{t}+\frac{b_\mu}{t^2},\qquad
a_\mu=\sqrt{\frac2{15}}\frac{\mu^2}{\pi},\quad b_\mu=\frac{\mu}{3\pi},\qquad t>0.
\]

The Jacobi normalization ensures \(|\lambda_j|\le N/2\). The exact half-filled sine-kernel boundary identity implies

\[
0\le M_{2,N}:=\operatorname{Tr}(J_N^2W_N)
\le\frac{32}{\pi^2}(H_N+1),
\]

and the refined exterior-lattice computation proves \(M_{2,N}=(3\pi^2)^{-1}\log N+O(1)\).

Define \(\eta(t)=\pi t\,{}_3F_2(\frac14,\frac12,\frac34;1,\frac32;4t^2)\) and \(V_N=-N[\eta(J_N/N)-\pi J_N/N]\). The positive-coefficient hypergeometric series gives

\[
\lvert\eta(t)-\pi t\rvert\le C_\eta|t|^3\quad(|t|\le\tfrac12),
\qquad C_\eta=4\pi\bigl[{}_3F_2(\tfrac14,\tfrac12,\tfrac34;1,\tfrac32;1)-1\bigr].
\]

All inputs are specific to the Wilson--Dirac/sine-kernel model and do not assume a reference-modular comparison.

\subsection{2. A fourth-moment bound for the positive majorant}\label{a-fourth-moment-bound-for-the-positive-majorant}

\textbf{Lemma 1.} For every even \(N\ge2\),

\[
\boxed{\operatorname{Tr}(G_NJ_N^4)=\sum_jq_j|\lambda_j|^4
\le\frac{a_\mu N^3}{6}+\frac{b_\mu N^2}{2}.}
\]

\textbf{Proof.} The diagonal representation holds because \(J_N^4\) is diagonal in the common reference eigenbasis, regardless of whether \(G_N\) commutes with \(J_N\). For the positive measure \(\nu_N=\sum_jq_j\delta_{|\lambda_j|}\), layer-cake integration and the Jacobi spectral radius give

\[
\begin{aligned}
\operatorname{Tr}(G_NJ_N^4)
&=4\int_0^{N/2}t^3\nu_N((t,\infty))\,dt\\
&\le4\int_0^{N/2}(a_\mu t^2+b_\mu t)\,dt
=\frac{a_\mu N^3}{6}+\frac{b_\mu N^2}{2}.
\end{aligned}
\]

The singular appearance of the spectral-tail bound near \(t=0\) causes no problem: the integrated majorant is finite. \(\square\)

\subsection{3. Improved transverse estimate}\label{improved-transverse-estimate}

\textbf{Theorem 2.} Let \(\mathcal Q_N^{\rm trans}=\operatorname{Tr}(D_NV_N)\). Then

\[
\boxed{
|\mathcal Q_N^{\rm trans}|
\le2C_\eta\sqrt{\left(\frac{a_\mu}{6N}+\frac{b_\mu}{2N^2}\right)M_{2,N}}
=O_\mu\!\left(\sqrt{\frac{\log N}{N}}\right).
}
\]

In particular \(\mathcal Q_N^{\rm trans}\to0\) as \(N\to\infty\) at each fixed \(\mu>0\).

\textbf{Proof.} In the reference eigenbasis, \(V_N\) has diagonal entries \(w_j\) with \(|w_j|\le C_\eta|\lambda_j|^3/N^2\). Using the pairing inequality followed by Cauchy--Schwarz with the \textbf{fourth positive-majorant moment}, not merely its total mass, yields

\[
\begin{aligned}
|\mathcal Q_N^{\rm trans}|
&\le\frac{2C_\eta}{N^2}\sum_j|\lambda_j|^3\sqrt{q_jb_j(1-b_j)}\\
&\le\frac{2C_\eta}{N^2}
\left(\sum_jq_j\lambda_j^4\right)^{1/2}
\left(\sum_j\lambda_j^2b_j(1-b_j)\right)^{1/2}\\
&=\frac{2C_\eta}{N^2}
\sqrt{\operatorname{Tr}(G_NJ_N^4)\operatorname{Tr}(J_N^2W_N)}.
\end{aligned}
\]

Apply Lemma 1 and the established Jacobi moment bound. \(\square\)

The finite-length estimate is explicit; the sharp asymptotic coefficient for \(M_{2,N}\) is not required to obtain the rate, only its \(O(\log N)\) upper bound.

\subsection{4. Bounded nonlinear correction and modular reduction}\label{bounded-nonlinear-correction-and-modular-reduction}

The established longitudinal estimate is

\[
|\mathcal Q_N^{\rm long}|\le C_\eta\left(\frac{3a_\mu}{8}+\frac{3b_\mu}{2N}\right).
\]

Therefore

\[
\boxed{
|\mathcal Q_N|\le C_\eta\left(\frac{3a_\mu}{8}+\frac{3b_\mu}{2N}\right)
+2C_\eta\sqrt{\left(\frac{a_\mu}{6N}+\frac{b_\mu}{2N^2}\right)M_{2,N}}
=O_\mu(1).}
\]

The exact two-bond asymptotic and finite quasifree relative-entropy decomposition are

\[
\mathcal P_L=\frac{\mu^2}{6}\log L+O_\mu(1),\qquad
\mathcal D_L=\mathcal P_L+\mathcal Q_N+\mathcal F_L-\Delta S_L,
\]

where \(\mathcal F_L=\operatorname{Tr}[(A_N-B_N)E_N]\) is the exact reference-modular residual. Theorem 1 of \emph{Uniform Square-Root-Mass Control of Wilson--Dirac Boundary Covariances}, with its entropy argument in section 5, proves \(\Delta S_L=O_\mu(1)\) for this Wilson scaling. Combining that imported bound with the nonlinear result of this paper yields

\[
\boxed{\mathcal D_L=\frac{\mu^2}{6}\log L+\mathcal F_L+O_\mu(1).}
\]

This reduction alone does not prove \(\mathcal F_L=o(\log L)\), convergence to a finite constant, or a proposed inverse-logarithmic convergence rate. The separate theorem in \emph{Uniform Bulk Modular Asymptotics and the Leading Wilson Entropy Law} proves \(\mathcal F_L=O_\mu(1)\), sufficient for the leading law. Theorem 3 of \emph{Uniform Modular Control and Nonlinear Convergence for Wilson--Dirac Intervals} proves \(\|E_N\|\le7\) for every even \(N\ge2\), under the exact concentration connection specified there. This later result is separate from the reduction proved above.

\subsection{5. Evidence and scope}\label{evidence-and-scope}

The fourth-moment estimate and theorem are exact mathematical consequences of the stated majorant and sine-kernel results. Independent finite arithmetic compares the improved bound with the earlier bound for eight Wilson configurations; all satisfy the stronger estimate. These tests are consistency checks, not the proof. The claimed rate holds at fixed \(\mu\) only, with no uniform assertion for \(\mu\) varying with \(N\).

\subsubsection{External mathematical context}\label{external-mathematical-context}

\begin{itemize}
\tightlist
\item
  V. Eisler and I. Peschel, ``Analytical results for the entanglement Hamiltonian of a free-fermion chain,'' \emph{J. Phys. A} \textbf{50} (2017) 284003, \url{doi:10.1088/1751-8121/aa76b5}.
\item
  D. Slepian, ``Prolate spheroidal wave functions, Fourier analysis, and uncertainty. V: The discrete case,'' \emph{Bell System Technical Journal} \textbf{57} (1978), 1371--1430, \url{doi:10.1002/j.1538-7305.1978.tb02104.x}.
\item
  S. Karnik, J. Romberg and M. A. Davenport, ``Improved bounds for the eigenvalues of prolate spheroidal wave functions and discrete prolate spheroidal sequences,'' \emph{Applied and Computational Harmonic Analysis} \textbf{55} (2021), 97--128, \url{doi:10.1016/j.acha.2021.04.002}.
\end{itemize}

These references support the classical concentration-operator setting; they do not independently prove the new Wilson-specific fourth-moment consequence.

\subsubsection{Acknowledgments and contributions}\label{acknowledgments-and-contributions}

Kirtus Justus directed the research. ChatGPT (OpenAI) developed and audited this fourth-moment derivation and its computational corroboration. Claude (Anthropic), Gemini (Google), Grok (SpaceXAI), and DeepSeek (High-Flyer) are acknowledged as broader research contributors without attribution of this derivation to them. There has been no external human peer review or institutional endorsement.

\end{document}
